\section{Conclusion}

The central claim of this paper is that the stability of SGD is not fully explained by pointwise properties of the loss landscape, such as local curvature. Instead, it also depends on temporal interactions between successive mini-batch updates, an aspect that remains comparatively underexplored despite its influence on the training trajectory. To address this gap, we introduced a signal--interference framework for stochastic optimization. By decomposing parameter growth into gradient signal and cross-step interference, the framework brings to light an important property of training dynamics: whether successive updates accumulate constructively or partially overwrite one another. We further derived a formulation to explain such interaction through consecutive-gradient cosine similarity, and used it to analyze plain SGD, momentum, and weight decay. Our results indicate that both stability and performance of SGD depend heavily on the combination of update magnitude and coherence. Momentum improves performance by making successive updates constructive, but it can become unstable if its velocity overshoots. Weight decay improves robustness by reducing the influence of older, stale directions. Overall, our findings complement existing approaches by introducing a temporal perspective on optimization dynamics, together with a practical diagnostic tool for assessing stability based on the correlation between successive updates. Therefore, these contributions provide a more mechanistic understanding of training evolution, helping to bridge the gap between theory and practice in modern machine learning.